% Options for packages loaded elsewhere
\PassOptionsToPackage{unicode}{hyperref}
\PassOptionsToPackage{hyphens}{url}
\PassOptionsToPackage{dvipsnames,svgnames,x11names}{xcolor}
\documentclass[
  10pt,
  fontset=fandol]{ctexart}
\usepackage{xcolor}
\usepackage[margin=16mm]{geometry}
\usepackage{amsmath,amssymb}
\setcounter{secnumdepth}{-\maxdimen} % remove section numbering
\usepackage{iftex}
\ifPDFTeX
  \usepackage[T1]{fontenc}
  \usepackage[utf8]{inputenc}
  \usepackage{textcomp} % provide euro and other symbols
\else % if luatex or xetex
  \usepackage{unicode-math} % this also loads fontspec
  \defaultfontfeatures{Scale=MatchLowercase}
  \defaultfontfeatures[\rmfamily]{Ligatures=TeX,Scale=1}
\fi
\usepackage{lmodern}
\ifPDFTeX\else
  % xetex/luatex font selection
\fi
% Use upquote if available, for straight quotes in verbatim environments
\IfFileExists{upquote.sty}{\usepackage{upquote}}{}
\IfFileExists{microtype.sty}{% use microtype if available
  \usepackage[]{microtype}
  \UseMicrotypeSet[protrusion]{basicmath} % disable protrusion for tt fonts
}{}
\makeatletter
\@ifundefined{KOMAClassName}{% if non-KOMA class
  \IfFileExists{parskip.sty}{%
    \usepackage{parskip}
  }{% else
    \setlength{\parindent}{0pt}
    \setlength{\parskip}{6pt plus 2pt minus 1pt}}
}{% if KOMA class
  \KOMAoptions{parskip=half}}
\makeatother
\setlength{\emergencystretch}{3em} % prevent overfull lines
\providecommand{\tightlist}{%
  \setlength{\itemsep}{0pt}\setlength{\parskip}{0pt}}
\usepackage{amsmath,amssymb,mathtools,needspace}
\xeCJKDeclareCharClass{CJK}{"2032,"0400->"04FF,"2160->"216B,"2460->"2473,"25A0,"25C7,"25CB,"25CF}
\IfFontExistsTF{Arial Unicode MS}{\xeCJKsetup{AutoFallBack=true}\setCJKfallbackfamilyfont{\CJKrmdefault}{Arial Unicode MS}}{}
\setcounter{secnumdepth}{0}
\setlength{\emergencystretch}{2em}
\allowdisplaybreaks
\usepackage{bookmark}
\IfFileExists{xurl.sty}{\usepackage{xurl}}{} % add URL line breaks if available
\urlstyle{same}
\hypersetup{
  pdftitle={Jacobi 迭代法},
  colorlinks=true,
  linkcolor={Maroon},
  filecolor={Maroon},
  citecolor={Blue},
  urlcolor={Blue},
  pdfcreator={LaTeX via pandoc}}

\title{Jacobi 迭代法}
\author{}
\date{}

\begin{document}
\maketitle

\subsubsection{8.2.1　Jacobi 迭代法}\label{jacobi-ux8fedux4ee3ux6cd5}

设有方程组

\[
\sum_{j=1}^n a_{ij}x_j=b_i\quad(i=1,2,\cdots,n),
\]

记作

\[
Ax=b,
\tag{8.2.1}
\]

\(A\) 为非奇异阵且 \(a_{ij}\neq0\)（\(i=1,2,\cdots,n\)）。将 \(A\) 分裂为 \(A=D-L-U\)，其中

\[
D=\begin{bmatrix}
a_{11}&&&&\\
&a_{22}&&&\\
&&\ddots&&\\
&&&\ddots&\\
&&&&a_{nn}
\end{bmatrix},\qquad
L=-\begin{bmatrix}
0&&&&\\
a_{21}&0&&&\\
a_{31}&a_{32}&0&&\\
\vdots&\vdots&\ddots&\ddots&\\
a_{n1}&a_{n2}&\cdots&a_{n,n-1}&0
\end{bmatrix},
\]

\[
U=-\begin{bmatrix}
0&a_{12}&a_{13}&\cdots&a_{1n}\\
&0&a_{23}&\cdots&a_{2n}\\
&&\ddots&\ddots&\vdots\\
&&&0&a_{n-1,n}\\
&&&&0
\end{bmatrix}.
\]

将式（8.2.1）第 \(i\)（\(i=1,2,\cdots,n\)）个方程用 \(a_{ii}\) 去除再移项，得到等价方程组

\begin{quote}
校注（agent 补充，ch08-E013）：本页式（8.2.1）下方原印 \(a_{ij}\neq0\)（\(i=1,2,\cdots,n\)），已按原样保留。这里没有给出 \(j\) 的范围；后续 Jacobi 分裂以对角元 \(a_{ii}\) 作除数，通常条件应为 \(a_{ii}\neq0\)。此处疑为原书下标错误；独立交叉复核指出后已回原图放大确认。
\end{quote}

\href{../../source-images/pdf-218.jpeg}{原页图像}

\[
x_i=\frac1{a_{ii}}\left(b_i-\sum_{\substack{j=1\\j\neq i}}^n a_{ij}x_j\right)\quad(i=1,2,\cdots,n),
\tag{8.2.2}
\]

简记作

\[
x=B_0x+f,
\]

其中

\[
B_0=I-D^{-1}A=D^{-1}(L+U),\qquad f=D^{-1}b.
\]

对方程组（8.2.2）应用迭代法，得到解式（8.2.1）的 \textbf{Jacobi 迭代公式}

\[
\begin{cases}
x^{(0)}=(x_1^{(0)},x_2^{(0)},\cdots,x_n^{(0)})^{\mathrm T}\quad\text{（初始向量）},\\[4pt]
x_i^{(k+1)}=\dfrac1{a_{ii}}\left(b_i-\displaystyle\sum_{\substack{i=1\\j\neq i}}^n a_{ij}x_j^{(k)}\right),
\end{cases}
\tag{8.2.3}
\]

其中 \(x^{(k)}=(x_1^{(k)},x_2^{(k)},\cdots,x_n^{(k)})^{\mathrm T}\) 为第 \(k\) 次迭代向量。设 \(x^{(k)}\) 已经算出，由式（8.2.3）可计算下一次迭代向量 \(x^{(k+1)}\)（\(k=0,\allowbreak 1,\allowbreak 2,\allowbreak \cdots;\ i=1,\allowbreak 2,\allowbreak \cdots,\allowbreak n\)）。

显然迭代公式（8.2.3）的矩阵形式为

\[
\begin{cases}
x^{(0)}\quad\text{（初始向量）},\\
x^{(k+1)}=B_0x^{(k)}+f,
\end{cases}
\tag{8.2.4}
\]

其中 \(B_0\) 称为 \textbf{Jacobi 方法迭代矩阵}。

由此看出，Jacobi 迭代法公式简单，每迭代一次只需计算一次矩阵和向量乘法。在用计算机计算时，需要两组工作单元，以存储 \(x^{(k)}\) 及 \(x^{(k+1)}\)。

\begin{center}\rule{0.5\linewidth}{0.5pt}\end{center}

\subsection{相关原页校注（agent 补充）}\label{ux76f8ux5173ux539fux9875ux6821ux6ce8agent-ux8865ux5145}

以下校注来自本节覆盖的原页，可能也涉及同页相邻小节；原文照录与校正说明分开保留。

\subsubsection{原书 PDF 页 218 的校注}\label{ux539fux4e66-pdf-ux9875-218-ux7684ux6821ux6ce8}

\href{../pages/pdf-218.md}{完整原页转录} · \href{../../source-images/pdf-218.jpeg}{原页图像}

校注记录：ch08-E002。

\begin{quote}
校注（agent 补充）：式（8.2.3）的求和下限原页明确印为 \(i=1\)，下一行仍为 \(j\neq i\)，已按原样保留。由本页式（8.2.2）可知，此处求和指标应为 \(j=1\)，属于疑似原书排印错误，并非 OCR 错误。
\end{quote}

\subsection{本节覆盖原页的校注记录}\label{ux672cux8282ux8986ux76d6ux539fux9875ux7684ux6821ux6ce8ux8bb0ux5f55}

下列记录按原页关联，含同页相邻小节的校注；计算时同时读取上文校注和完整原页。

\begin{itemize}
\tightlist
\item
  \texttt{ch08-E013}：原书 PDF 页 217
\item
  \texttt{ch08-E002}：原书 PDF 页 218
\end{itemize}

\end{document}
